\section{Tetrahedral null-generated causal cone}
Define
\[
\mathbf n_1=\frac{(1,1,1)}{\sqrt3},\quad
\mathbf n_2=\frac{(1,-1,-1)}{\sqrt3},\quad
\mathbf n_3=\frac{(-1,1,-1)}{\sqrt3},\quad
\mathbf n_4=\frac{(-1,-1,1)}{\sqrt3}.
\]
Then
\[
\sum_a\mathbf n_a=0,
\qquad
\mathbf n_a\cdot\mathbf n_b=-\frac13\ (a\ne b),
\qquad
\frac14\sum_a\mathbf n_a\mathbf n_a^T=\frac13I_3.
\]
With Minkowski metric $\eta=\diagop(+,-,-,-)$ define $k_a=(1,\mathbf n_a)$. Each $k_a$ is future null, and
\[
\frac14\sum_a k_a=(1,0,0,0).
\]

For $p_a\ge0$, set $P=\sum_ap_ak_a$.
\begin{theorem}[Exact tetrahedral causal norm]\label{thm:cone}
\[
\boxed{P^2=\frac83\sum_{a<b}p_ap_b\ge0.}
\]
Thus the positive span of the four null rays is contained in the future Lorentz cone, and $P$ is null exactly when at most one $p_a$ is nonzero.
\end{theorem}
\begin{proof}
$k_a^2=0$, while for $a\ne b$,
\[
k_a\cdot k_b=1-\mathbf n_a\cdot\mathbf n_b=\frac43.
\]
Expanding $P^2$ gives the formula.
\end{proof}
At fixed $P^0=1$ the section is the regular tetrahedron inside the unit ball. Its four vertices are null directions; nontrivial convex mixtures are timelike. It is therefore a null-generated tetrahedral causal cone, not a replacement for the full null boundary.

\section{Spinorial and supersymmetric lift}
Let $\sigma^\mu=(I_2,\sigma_x,\sigma_y,\sigma_z)$ and
\[
K_a=I_2+\mathbf n_a\cdot\boldsymbol\sigma.
\]
Since $|\mathbf n_a|=1$, $K_a$ is positive semidefinite, rank one, and $\det K_a=0$. Hence there is a two-spinor $\lambda_a\in\C^2$ such that
\[
\boxed{K_a=\lambda_a\lambda_a^\dagger}.
\]
The factorization has the full phase gauge $\lambda_a\sim e^{i\chi_a}\lambda_a$; the sign $\lambda_a\mapsto-\lambda_a$ is the Euler half-turn element, not the entire gauge freedom.

The geometry supplies the null spinors but not fermionic operators. Attach four modes $f_a,f_a^\dagger$ satisfying
\[
\{f_a,f_b^\dagger\}=\delta_{ab},
\qquad
\{f_a,f_b\}=\{f_a^\dagger,f_b^\dagger\}=0.
\]
For $p_a\ge0$ define
\[
P_{\alpha\dot\beta}=\sum_a p_a\lambda_{a\alpha}\bar\lambda_{a\dot\beta},
\]
\[
Q_\alpha=\sqrt2\sum_a\sqrt{p_a}\lambda_{a\alpha}f_a^\dagger,
\qquad
\bar Q_{\dot\beta}=\sqrt2\sum_a\sqrt{p_a}\bar\lambda_{a\dot\beta}f_a.
\]

\begin{theorem}[Tetrahedral $\mathcal N=1$ supertranslation algebra]\label{thm:tetra-susy}
\[
\boxed{\{Q_\alpha,\bar Q_{\dot\beta}\}=2P_{\alpha\dot\beta}},
\qquad
\{Q_\alpha,Q_\beta\}=0,
\qquad
\{\bar Q_{\dot\alpha},\bar Q_{\dot\beta}\}=0.
\]
Thus the tetrahedral momentum cone admits an exact $3+1$ dimensional $\mathcal N=1$ supertranslation realization once one CAR mode is supplied for each null ray.
\end{theorem}
\begin{proof}
Using the CAR,
\[
\{Q_\alpha,\bar Q_{\dot\beta}\}
=2\sum_{a,b}\sqrt{p_ap_b}\lambda_{a\alpha}\bar\lambda_{b\dot\beta}\{f_a^\dagger,f_b\}
=2\sum_ap_a\lambda_{a\alpha}\bar\lambda_{a\dot\beta}.
\]
The same-chirality anticommutators vanish because the creation operators mutually anticommute, as do the annihilation operators.
\end{proof}

Theorem~\ref{thm:tetra-susy} proves that the tetrahedral spinorial carrier admits a natural supersymmetric lift after CAR input. It does not prove that geometry alone derives physical fermions.
